\selectlanguage{ngerman}%
\begin{ledgroupsized}[r]{120mm}%
\footnotesize%
\pstart%
\noindent\textbf{Überlieferung:}%
\pend%
\end{ledgroupsized}%
\begin{ledgroupsized}[r]{114mm}%
\footnotesize%
\pstart%
\parindent -6mm%
\makebox[6mm][l]{\textit{L}}%
Aufzeichnung:
LH~XXXVII~5 Bl.\;39. 
Ein Blatt~8\textsuperscript{o},
Ränder beschnitten;
Papiererhaltungsmaßnahmen.
Zwei Seiten.
\pend%
\end{ledgroupsized}%
\vspace{5mm}%
\begin{ledgroup}%
\footnotesize%
\pstart%
\noindent%
\textbf{Datierungsgründe:}
Die Datierung der Aufzeichnung N.\;\ref{60295} beruht %
auf inhaltlichen Betrachtungen.
Siehe die ausführliche Begründung auf S.\;\pageref{60295_dat}.
\pend%
\end{ledgroup}%
\selectlanguage{latin}%
\frenchspacing%
\vspace{8mm}
\pstart%
\normalsize%
\noindent%
\lbrack%
\href{\Scan{LH_37_05_039r}}%
{39\,r\textsuperscript{o}}%
\rbrack\ %
Si punctum
\edtext{%
moveatur motu composito ex duobus,%
\indexS{motus compositus}%
}{%
\lemma{moveatur}%
\Bfootnote{%
\textit{(1)}~duobus motibus, %
\textit{(2)}~motu composito ex duobus,%
~\textit{L}%
}}
quorum
\edtext{%
quilibet habeat celeritates
aequalibus spatii elementis%
\indexS{elementum spatii}
eadem proportione decrescentes%
\indexS{celeritas decrescens}%
\lbrack,\rbrack%
}{%
\lemma{quilibet}%
\Bfootnote{%
\textit{(1)}~decrescat %
\textit{(2)}~habeat celeritates %
\lbrack...\rbrack\ proportione decrescentes%
~\textit{L}%
}}
adeoque in
\edtext{%
progressione Geometrica,%
\indexS{progressio geometrica}%
}{%
\lemma{progressione}%
\Bfootnote{%
\textit{(1)}~Arithmetica,\indexS{progressio arithmetica} %
\textit{(2)}~Geometrica,%
~\textit{L}%
}}
videndum
an motus compositus%
\indexS{motus compositus}
sit etiam in progressione Geometrica.%
\indexS{progressio geometrica}
\pend%
\pstart%
Videndum etiam
quid fiat in casu tali%
\lbrack:\rbrack\
si corpus \textit{A} simul incidat%
\indexS{corpus incidens in alia plura}
in corpus \textit{B} %
\edtext{%
et \textit{C},
tunc ob incursum%
}{%
\lemma{et \textit{C},}%
\Bfootnote{%
\textit{(1)}~ob incursum %
\textit{(2)}~tunc ob incursum%
~\textit{L}%
}}
in \textit{B} %
retardabitur ejus velocitas progrediendi%
\indexS{velocitas progrediendi}%
\indexS{velocitas retardata}
in recta \textit{EB} et parallelis
in
\edtext{%
ratione $A - B$ ad $A + B;$%
\indexS{ratio retardationis}%
}{%
\lemma{ratione}%
\Bfootnote{%
\textit{(1)}~$A + B$ %
\textit{(2)}~$A - B$ ad $A + B;$%
~\textit{L}%
}}
sed ejusdem velocitas progrediendi%
\indexS{velocitas progrediendi}
in recta \textit{CD} et parallelis %
diminuetur in ratione $A - C$ ad $A + B;$
quomodo ergo has diminutiones componemus inter se,%
\indexS{diminutio velocitatis}%
\indexS{compositio diminutionum velocitatis}
et quem motum dabimus toti.%
\indexS{motus totius}
Neque enim fertur corpus \textit{A}
motu composito in \textit{CD} %
et in \textit{EB},
sed motu in \textit{EB} %
qui includit motum in \textit{CD}.%
\indexS{motus compositus}
An
\edtext{%
ergo resolvemus motum
in \textit{EB} in duos,%
}{%
\lemma{ergo}%
\Bfootnote{%
\textit{(1)}~componemus mot %
\textit{(2)}~resolvemus motum in %
\textit{(a)}~\textit{AB} %
\textit{(b)}~\textit{EB} in duos,%
~\textit{L}%
}}
unum in \textit{EF} seu \textit{DC},
alterum in \textit{FB};
ex quibus motus in \textit{FB} retardabitur
in ratione $A - B$ ad $A + B.$%
\indexS{motus retardatus}%
\indexS{ratio retardationis}%
\pend%
\vspace{1.0em} %
\centerline{%
\includegraphics[width=0.28\textwidth]{%
gesamttex/edit_VIII-4/images/LH_37_05_039_d1_039r.pdf%
}}
\nobreak\vspace{0.0em}%
\centerline{%
\lbrack\textit{Fig.\:1}\rbrack%
}
\pstart%
\noindent%
\edlabel{LH_37_05_039_039r_xjh-1}%
\edtext{}{%
\xxref{LH_37_05_039_039r_xjh-1}{LH_37_05_039_039r_xjh-2}%
\lemma{Sed}%
\Bfootnote{%
\textit{(1)}~motus in \textit{DC} vel \textit{EF} %
\textit{(a)}~quomodo retardabitur %
\textit{(b)}~si etiam in certa quadam ratione %
\textlangle\textendash\textrangle\ retardetur, hinc jam redibit quaestio\indexS{quaestio} initio posita, haud dubie affirmative resolvenda. Sed quaeritur utrum simul %
\textit{(2)}~quaeritur utrum %
\lbrack...\rbrack\ simul in ratione %
\textit{(a)}~$A + B$ %
\textit{(b)}~$A - B$ ad $A + B$ ob %
\textit{(aa)}~\textit{A} %
\textit{(bb)}~\textit{B} et in ratione $A - C$ ad $A + C$ %
\bars{motus in \textit{CD} vel \textit{DF} retardetur \textit{erg. u. gestr.}}%
, ita ut simul retroagatur differentia %
\textit{(aaa)}~inter %
\textit{(bbb)}~quae est %
\lbrack...\rbrack\ quae est %
\bars{in \textit{ändert Hrsg.}} celeritatem $\overline{A + B} : \overline{A + B}$ et $\overline{A - B}$ ad $\overline{A + B}.$%
~\textit{L}%
}}%
Sed quaeritur
utrum retardetur motus in \textit{DC} vel \textit{EF}
in ratione $A - B - C$ ad $A + B + C;$%
\indexS{ratio retardationis}
an vero simul in ratione $A - B$ ad $A + B$
ob \textit{B} %
et in ratione $A - C$ ad $A + C$,%
\indexS{ratio retardationis}
ita ut simul retroagatur differentia%
\indexS{differentia celeritatum}
quae est inter celeritatem
$\overline{A + C} : \overline{A + C}$
et
$\overline{A - C} : \overline{A + C}$
et quae est
in\lbrack ter\rbrack\
celeritatem
$\overline{A + B} : \overline{A + B}$
et
$\overline{A - B}$ ad $\overline{A + B}.$%
\edlabel{LH_37_05_039_039r_xjh-2}
Differentiae%
\indexS{differentia celeritatum}
autem sunt
$2C : \overline{A + C}$
et
$2B : \overline{A + B},$
quae in unum additae
faciunt
$2C \cdot \overline{A + B} + 2B \cdot \overline{A + C}\, , \, :\overline{\overline{A + C} \cdot \overline{A + B}}$ %
quae est celeritas retroactionis%
\indexS{celeritas retroactionis}%
\indexS{retroactio}%
\lbrack,\rbrack\
quae detracta a celeritate prima
$\overline{A + C} \cdot \overline{A + B} : \overline{A + C} \cdot \overline{A + B},$
seu 1,
dabit celeritatem residuam%
\indexS{celeritas residua}
$AA + AB + AC + BC - 2AC - 2BC - 2AB - 2BC\, , \, : \overline{A + B} \cdot \overline{A + C}$
seu
$\overline{AA - AB - AC - 3BC} : \overline{AA + AB + AC + BC}.$
Unde alius prodit calculus.%
\indexS{calculus}
Utrum igitur eligi
\edlabel{LH_37_05_039_039rv_gdh-1}%
debet?%
\edtext{}{%
\xxref{LH_37_05_039_039rv_gdh-1}{LH_37_05_039_039rv_gdh-2}%
\lemma{debet?}%
\Bfootnote{%
\textit{(1)}~Sit %
\textit{(2)}~Et sane. \lbrack/\rbrack\ Si ponamus%
~\textit{L}%
}}
Et sane.
\lbrack%
\href{\Scan{LH_37_05_039v}}%
{39\,v\textsuperscript{o}}%
\rbrack%
\pend%
\pstart%
Si ponamus%
\edlabel{LH_37_05_039_039rv_gdh-2}
\edtext{%
corpus \textit{G}%
}{%
\lemma{corpus \textit{G}}%
\Cfootnote{Vgl. das Diagramm \lbrack\textit{Fig.\:2}\rbrack.%
}}
simul incurrere
in corpora quiescentia \textit{H} et \textit{L},%
\indexS{corpus incurrens in alia quiescentia}%
\indexS{corpus quiescens}
perinde res prodibit
ac si \textit{H} et \textit{L} essent unum
\edtext{%
corpus,
seu retardabitur velocitas%
\indexS{velocitas retardata}
in ratione $G - H - L$ ad $G + H + L.$
Itaque non simul retardabitur
in ratione $G - H$ ad $G + H$
et in ratione $G - L$ ad $G + L.$%
\indexS{ratio retardationis}
Differentem enim%
}{%
\lemma{corpus,}%
\Bfootnote{%
\textit{(1)}~nec di %
\textit{(2)}~seu retardabitur velocitas in ratione %
\textit{(a)}~$H - L$ %
\textit{(b)}~\textit{G} %
\textit{(c)}~$G - H - L$ ad $G + H + L.$ %
\textit{(aa)}~Non vero %
\textit{(bb)}~Itaque non simul %
\textit{(aaa)}~in ratione %
\textit{(bbb)}~retardabitur in %
\lbrack...\rbrack\ Differentem enim%
~\textit{L}%
}}
rationem prodituram%
\indexS{ratio retardationis}
calculus praecedens ostendit.%
\indexS{calculus}
Hinc patet
quam facile sit
compositionibus illis abuti%
\indexS{compositio motuum}%
\lbrack,\rbrack\
praesertim %
cum si
\edtext{%
quemlibet incursum%
\indexS{incursus plurimus}
separatim consideres%
\lbrack,\rbrack%
}{%
\lemma{quemlibet}%
\Bfootnote{%
\textit{(1)}~separatim consideres %
\textit{(2)}~incursum separatim consideres%
~\textit{L}}}
is det impulsum Eccentricum.%
\indexS{impulsus eccentricus}
\pend%
\pstart%
Licebit tamen et incursus duos%
\indexS{incursus plurimus}
simul comminisci %
quorum neuter separatim det impulsum Eccentricum,%
\indexS{impulsus eccentricus}
ut ita haec quoque excusatio cesset,%
\indexS{excusatio}
ut si
\edtext{%
corpus \textit{M}%
}{%
\lemma{corpus \textit{M}}%
\Cfootnote{%
Vgl. das Diagramm \lbrack\textit{Fig.\:3}\rbrack.%
}}
simul incurrat
in corpora \textit{N} et \textit{O}.%
\indexS{corpus incurrens in alia plura}
\edlabel{LH_37_05_039_039v_principium-1}%
Hinc patet
generalibus principiis%
\indexS{principium generale}
insistendum %
de manente vi respectiva unius corporis%
\indexS{vis respectiva unius corporis}
secundum reliqua omnia simul sumta%
\lbrack,\rbrack\
sumendo ipsa pro
\edlabel{LH_37_05_039_039v_ykf-1}%
\edtext{}{%
\xxref{LH_37_05_039_039v_ykf-1}{LH_37_05_039_039v_ykf-2}%
\lemma{uno. \lbrack/\rbrack}\Bfootnote{%
\textit{(1)}~Posito igitur %
\textit{(2)}~In prima figura %
\textit{(3)}~Videndum igitur %
\lbrack...\rbrack\ figura sic %
\textit{(a)}~progredi %
\textit{(b)}~procedere liceat,%
~\textit{L}%
}}%
uno.%
\edlabel{LH_37_05_039_039v_principium-2}%
\pend%
\pstart%
Videndum igitur porro
an in
\edtext{%
prima figura%
\indexS{figura}%
}{%
\lemma{prima figura}%
\Cfootnote{Vgl. das Diagramm \lbrack \textit{Fig.\:1}\rbrack.%
}}
sic procedere liceat,%
\edlabel{LH_37_05_039_039v_ykf-2}
ut celeritas ipsius \textit{A}%
\indexS{celeritas corporis incurrentis}
secundum \textit{EB} %
retardetur quidem in ratione $A - B$ ad $A + B,$%
\indexS{celeritas retardata}%
\indexS{ratio retardationis}
celeritas autem secundum \textit{EF} vel \textit{DC} %
in ratione $A - B - C$%
\indexS{celeritas retardata}%
\indexS{ratio retardationis}
ad
\edtext{%
$A + B + C;$
sed hoc facto%
\lbrack,\rbrack\
si prius celeritates compositae%
\indexS{celeritas composita}
in \textit{EF} et \textit{FB} %
dedere \textit{EB},%
}{%
\lemma{$A + B + C;$}%
\Bfootnote{%
\textit{(1)}~quo facto haud dubie celeritas inde composita in ratione etiam certa in \textit{AB}, etiam in ratione certa retardabitur %
\textit{(2)}~sed hoc facto %
\textit{(a)}~celerit %
\textit{(b)}~si prius celeritates compositae, %
\textit{(aa)}~dedere %
\textit{(bb)}~in \textit{EF} et \textit{FB} dedere \textit{EB},%
~\textit{L}%
}}
nunc mutata proportione%
\indexS{proportio celeritatum}
dabunt aliam lineam;
\edtext{%
si ex \textit{EF} fiat \textit{EQ}%
}{%
\lemma{si ex \textit{EF} fiat \textit{EQ}}%
\Cfootnote{Vgl. das Diagramm \lbrack\textit{Fig.\:4}\rbrack, welches auf das Diagramm \lbrack\textit{Fig.\:1}\rbrack\ bezogen ist.
Der Buchstabe \textit{Q} bezeichnet in \lbrack\textit{Fig.\:4}\rbrack\ aber nicht denselben Punkt bzw. Körper wie in \lbrack\textit{Fig.\:1}\rbrack.%
}}
et ex \textit{FB} fiat
\edtext{%
\textit{QR}.
At%
}{%
\lemma{\textit{QR}}%
\Bfootnote{%
\textit{(1)}~, nec %
\textit{(2)}~. At%
~\textit{L}%
}}
licet simul incurratur in corpora%
\indexS{incursus plurimus}
\edtext{%
\edtext{%
\textit{B}, \textit{C} et
\textit{Q},%
}{%
\lemma{corpora \textit{B}, \textit{C} et \textit{Q}}%
\Cfootnote{Vgl. das Diagramm \lbrack\textit{Fig.\:1}\rbrack.}}
ubi servanda est linea \textit{EB},
videndum
an servetur eadem quoque semper proportio ipsius \textit{ER}.%
\indexS{proportio celeritatum}%
}{%
\lemma{\textit{B}, \textit{C} et \textit{Q},}%
\Bfootnote{%
\textit{(1)}~ita ut %
\textit{(a)}~simul %
\textit{(b)}~eadem servetur linea %
\textit{(2)}~ubi servanda est linea \textit{EB}, %
\textit{(a)}~videtur %
\textit{(b)}~videndum an %
\lbrack...\rbrack\ ipsius \textit{ER}.%
~\textit{L}%
}}
\pend%
\pstart%
Verum non est data haec compositionum Methodus.%
\indexS{methodus compositionis motuum}%
\indexS{compositio motuum}
Sed recurrendum est ad
\edtext{%
principia mea generalia,%
\indexS{principium generale}%
\indexS{principium meum}
\edlabel{LH_37_05_039_039v_principium-3}%
quod eadem semper servetur vis respectiva%
\indexS{vis respectiva servanda}
inter unum corpus et reliqua plura simul sumta;%
\edlabel{LH_37_05_039_039v_principium-4}%
}{%
\lemma{principia \lbrack...\rbrack\ sumta}%
\Cfootnote{%
Siehe zum methodischen Prinzip der Erhaltung der \textit{potentia} bzw. \textit{vis respectiva}
etwa G.\:W.\;\textsc{Leibniz},
\textit{Dynamica}, pars II, sectio III, prop.\;7 bis prop.\;12
(%
\textit{CFP}~I, S.\;400--412~\bracks{A};
II.2, S.\;498--512~\bracks{C})\cite{01354}
sowie die Aufzeichnungen
N.\;\ref{60111},
N.\;\ref{58220},
N.\;\ref{58235},
N.\;\ref{58949}, %
und
N.\;\ref{60298} %
im vorliegenden Band.
Siehe zudem die ausführlichen Datierungsgründe zu N.\;\ref{60295}, unten S.\;\pageref{60295_dat}.%
}}
itaque cum absit gyratio%
\indexS{gyratio in concursu eccentrico}
seu impulsus eccentricus%
\indexS{impulsus eccentricus}
hoc loco,
cum in
\edtext{%
\textit{C}, \textit{B}, \textit{Q}%
}{%
\lemma{\textit{C}, \textit{B}, \textit{Q}}%
\Cfootnote{%
Vgl. das Diagramm \lbrack\textit{Fig.\:1}\rbrack.%
}}
simul incurritur,%
\indexS{incursus plurimus}
utique sumtis ipsis pro uno
simul retardabitur celeritas ipsius \textit{A}%
\indexS{celeritas retardata}
in ratione%
\indexS{ratio retardationis}
\edtext{%
$A - B - C - [Q]$
ad
$A + B + C + [Q],$
et%
}{%
\lemma{}%
\Bfootnote{%
$A - B - C - \bars{D \ \text{\textit{ändert Hrsg.}}}$%
ad
$A + B + C + \bars{D \ \text{\textit{ändert Hrsg.}}},$
et%
\textit{L}%
}}
hoc scilicet
quando corpus inter movendum
\textlangle non\textrangle\
\edtext{%
gyratur,
et semper \textit{EB}%
}{%
\lemma{gyratur,}%
\Bfootnote{%
\textit{(1)}~\textit{AB} %
\textit{(2)}~et semper \textit{EB}%
~\textit{L}%
}}
est
\edtext{\lbrack pars\rbrack%
}{%
\lemma{}\Bfootnote{%
\textlangle tan\textrangle gens
\textit{L~ändert Hrsg.}}}
lineae
quae describitur.%
\indexS{linea descripta}
\pend%
\vspace{3.0em} %
\centerline{\hspace{-95mm}%
\includegraphics[width=0.17\textwidth]{%
gesamttex/edit_VIII-4/images/LH_37_05_039_d2_039v.pdf%
}}
\nobreak\vspace{0.5em}%
\centerline{\hspace{-95mm}%
\lbrack\textit{Fig.\:2}\rbrack%
}
\vspace{-8.625em} %
\centerline{%
\includegraphics[width=0.30\textwidth]{%
gesamttex/edit_VIII-4/images/LH_37_05_039_d3_039v.pdf%
}}
\nobreak\vspace{0.5em}%
\centerline{%
\lbrack\textit{Fig.\:3}\rbrack%
}
\vspace{-8.35em} %
\centerline{\hspace{90mm}%
\includegraphics[width=0.145\textwidth]{%
gesamttex/edit_VIII-4/images/LH_37_05_039_d4_039v.pdf%
}}
\nobreak\vspace{0.65em}%
\centerline{\hspace{90mm}%
\lbrack\textit{Fig.\:4}\rbrack%
}
\newpage%
\pstart%
\noindent%
\centering%
*    *    *
\pend%
\vspace{3.0em}%
\selectlanguage{ngerman}%
\begin{ledgroup}%
\footnotesize%
\pstart%
\noindent%
\textls{Ausführliche Datierungsgründe:}
\pagelabel{60295_dat}%
Bei der titellosen und undatierten Aufzeichnung N.\;\ref{60295} bieten die Eigenschaften des Textträgers keine Anhaltspunkte für eine Datierung dar.
Aus inhaltlichen Gründen erweist sich eine grobe chronologische Einordnung des Textes dennoch als möglich.
Im Mittelpunkt der Aufzeichnung steht die Kinematik des mehrfachen Stoßes, insbesondere die Frage nach der Verzögerung eines Körpers, der zwei oder mehrere Körper zugleich anstoßt.
Zur Beantwortung dieser Frage nimmt Leibniz im Text ausdrücklich einen methodischen Grundsatz in Anspruch, dem zufolge in einem mehrere bewegte Körper umfassenden System \textit{eadem semper servetur vis respectiva inter unum corpus et reliqua plura simul sumta} (\refpassage{LH_37_05_039_039v_principium-3}{LH_37_05_039_039v_principium-4}; vgl. auch \refpassage{LH_37_05_039_039v_principium-1}{LH_37_05_039_039v_principium-2}).
Dieser methodische Grundsatz %
ist bereits in dem zwischen Rom\indexO{Rom} und Florenz\indexO{Florenz} verfassten Teil der \textit{Dynamica} ein wichtiger Bestandteil der Stoßlehre (vgl. pars II, sectio III, %
prop.\;7 bis prop.\;12
(%
\textit{CFP}~I, S.\;400--412~\bracks{A};
II.2, S.\;498--512~\bracks{C})\cite{01354}
und wird in Texten der späten 1680er und frühen 1690er Jahre mehrfach erläutert und angewendet;
vgl. etwa N.\;\ref{60111}, N.\;\ref{58220}, N.\;\ref{58235}, N.\;\ref{58949}, N.\;\ref{41780}, N.\;\ref{58947}, N.\;\ref{60298} und N.\;\ref{58293} im vorliegenden Band.
Auch die Aufzeichnung N.\;\ref{60295} ist vermutlich im gleichen Zeitraum entstanden.
Insbesondere könnte sie -- wie die Aufzeichnung N.\;\ref{60339} -- zwischen dem Textkomplex N.\;\ref{sdie_intro} \textit{Schedae de impulsu eccentrico} (von Leibniz eigenhändig auf den 31.\;Oktober\,/\,10.\;November 1690 datiert) und dem Textkomplex N.\;\ref{cecs_intro} \textit{Concursus eccentrici casus simplicissimus} (von Leibniz eigenhändig auf September 1691 datiert) verfasst worden sein.
Eine frühere Entstehung von N.\;\ref{60295} ist nicht auszuschlieschen; es ist aber unwahrscheinlich, dass die Aufzeichnung auf die Zeit vor Leibnizens Rückkehr aus der Italienreise\indexO{Italien} zurückgeht.
Eine spätere Entstehung von N.\;\ref{60295} kann ebensowenig ausgeschlossen werden.
\pend%
\end{ledgroup}%
\normalsize%
\selectlanguage{latin}%
\frenchspacing%
